\documentclass[a4paper,12pt,twoside,openright,titlepage]{book}

%Additional packages
\usepackage[utf8]{inputenc}
\usepackage[T1]{fontenc}
\usepackage[dutch,english]{babel}
\usepackage{imakeidx}
\usepackage{syntonly}
\usepackage[official]{eurosym}
%\usepackage[graphicx]
\usepackage{graphicx}
\graphicspath{ {./images/} }
\usepackage{float}
\usepackage{xurl}
\usepackage{hyperref}
\hypersetup{colorlinks=true, linkcolor=blue, citecolor=blue, filecolor=blue, urlcolor=blue, pdftitle=, pdfauthor=, pdfsubject=, pdfkeywords=}
\usepackage{tabularx}
\usepackage[table]{xcolor} % Table colors
\usepackage{scrextend}
\addtokomafont{labelinglabel}{\sffamily}
\usepackage{listings}
\usepackage{adjustbox}
\usepackage{color}

% Define colors
\definecolor{ashgrey}{rgb}{0.7, 0.75, 0.71}

% Listing style
\lstset{
  backgroundcolor=\color{ashgrey}, % choose the background color; you must add \usepackage{color} or \usepackage{xcolor}; should come as last argument
  basicstyle=\footnotesize,        % the size of the fonts that are used for the code
  breakatwhitespace=true,          % sets if automatic breaks should only happen at whitespace
  breaklines=true,                 % sets automatic line breaking
  extendedchars=true,              % lets you use non-ASCII characters; for 8-bits encodings only, does not work with UTF-8
  frame=single,	                   % adds a frame around the code
  rulecolor=\color{black},         % if not set, the frame-color may be changed on line-breaks within not-black text (e.g. comments (green here))
  keepspaces=true,                 % keeps spaces in text, useful for keeping indentation of code (possibly needs columns=flexible)
  columns=fullflexible,		   % make copy and paste possible
  showstringspaces=false,          % if true show spaces in strings adding particular underscores
  showspaces=false,                % if true show spaces everywhere adding particular underscores; it does not override 'showstringspaces'
}

% Uncomment for production
% \syntaxonly

% Style
\pagestyle{headings}

% Turn on indexing
\makeindex[intoc]

% Define document
\author{D. Leeuw}
\title{Trouble Shooting}
\date{\today\\v.0.8.0}

\begin{document}
\selectlanguage{dutch}

\maketitle

\copyright\ 2023 Dennis Leeuw\\

\begin{figure}
\includegraphics[width=0.3\textwidth]{CC-BY-SA-NC.png}
\end{figure}

\bigskip

Dit werk is uitgegeven onder de Creative Commons BY-NC-SA Licentie en laat anderen toe het werk te kopi\"eren, distribueren, vertonen, op te voeren, en om afgeleid materiaal te maken, zolang de auteurs en uitgever worden vermeld als maker van het werk, het werk niet commercieel gebruikt wordt en afgeleide werken onder identieke voorwaarden worden verspreid.


%%%%%%%%%%%%%%%%%%%
%%% Introductie %%%
%%%%%%%%%%%%%%%%%%%

\frontmatter
%\chapter{Over dit Document}
%\input{src/OverDitDocument}
%\input{src/DocChanges}

%%%%%%%%%%%%%%%%%
%%% De inhoud %%%
%%%%%%%%%%%%%%%%%
\tableofcontents

\mainmatter
\chapter{Inleiding}
Dit document neemt je mee in het proces van probleem oplossen, of in het Engels Troublshooting. Het behandelt welke stappen je kunt doorlopen en hoe je ervoor zorgt dat je tot een oplossing komt.


\chapter{Aannemen}
Bij elk probleem:
\begin{itemize}
	\item Noteer datum en tijd
	\item Noteer de naam van de contactpersoon (je eigennaam als je het zelf bent)
	\item Maak notities!
\end{itemize}

Als je zelf een probleem hebt ontdekt:
\begin{itemize}
\item Maak notities
	\begin{itemize}
	\item Als het mogelijk is: Maak een kopie van de error melding (copy/past, screenshot)
	\item Is dit een nieuwe omgeving of heeft het voorheen goed gewerkt?
	\item Als het voorheen goed gewerkt heeft, wat is er veranderd?
	\item Noteer de omstandigheden (wat was je aan het doen, hoe viel het je op)
	\item Noteer andere zaken die met het probleem te maken kunnen hebben (uitvallen lampjes, beeldscherm, etc.)
	\item Is het probleem herhaalbaar?
	\end{itemize}
\end{itemize}

Een probleem kan gebracht worden door iemand anders. In dat geval:
\begin{itemize}
	\item Maak notities
	\item Vraag door bij onduidelijkheden, controleer de voorgaande lijst om te zien of je alle relevante informatie hebt.
	\item Als de persoon klaar is met zijn verhaal, vat het verhaal samen en controlleer of je goed begrepen hebt wat het probleem is.
	\item Vraag hoe je contact kan houden (e-mail, telefoon, face-to-face?)
	\item Vraag of er meerdere mensen met hetzelfde probleem zijn?
\end{itemize}

Documenteer het probleem in een ticketsysteem. Een centraalsysteem om problemen in op te slaan heeft voordelen.
\begin{itemize}
	\item Anderen kunnen je probleem ook zien (helpen)
	\item Anderen kunnen een probleem van een gebruiker/klant overnemen als je ziek bent
	\item Veel voorkomende problemen worden sneller duidelijk een fix kan gedocumenteerd worden
\end{itemize}


\chapter{Ademhalen}
Doe niets! Haal adem! Een ontspannen geest zorgt voor een beter denkproces. Denk nog even na over het probleem.



\chapter{Analyse}
\begin{itemize}
\item Zijn er anderen met hetzelfde probleem
\item Is er een oplossing die je kan gebruiken (knowledge base)
\item Zijn er wijzigingen doorgevoerd die het probleem kunnen verklaren?
	\begin{itemize}
	\item Patches: Microsoft patch tuesday
	\item Nieuw programma
	\item Nieuwe hardware
	\item Nieuwe gebruiker
	\end{itemize}
\item Hardware issues:
	\begin{itemize}
		\item kabel los
		\item aan/uit schakelaar
		\item Oude hardware
	\end{itemize}
\item Documenteer wat je hebt ontdekt
\end{itemize}


\chapter{Onderzoek}
\begin{enumerate}
\item Kijk in de helpdesk applicatie of het probleem eerder gemeld is en of er een oplossing voor is
\item Zoek de letterlijke error-melding in een zoekmachine
\item Je bent waarschijnlijk niet de eerste met het probleem, gebruik een zoekmachine
\item Gebruik andere woorden in de zoekmachine
\item Vraag een collega om hulp
\item Documenteer de gevonden oplossingen en de uitkomsten van je tests
\end{enumerate}

\section{Zorg voor jezelf}
Vergeet niet om voor jezelf te zorgen. De beste oplossingen komen vaak op de WC. We nemen tijdens het storingszoeken vaak te weinig pauze en als we dan gedwongen worden door onze blaas om naar de WC te gaan komt soms spontaan de oplossing. Dit is een les. We moeten vaker een pauze nemen om helder te kunnen denken.

Als je uren achter elkaar achter een probleem aan hebt gezeten voel je dan niet schuldig om een kwartier pauze te nemen. Ga wandelen, zoek de frisse lucht op, zorg dat je even weg bent van je werkplek (zelfs als er een server down is of het halve bedrijf niet kan werken). Even rust zorgt vaak voor de oplossing en zorgt er zeker voor dat je daarna helderder nadenkt.

Vergeet ook de interne mens niet. Er is niets zo chagrijning makend als honger of dorst. Sla dus geen maaltijden over, denk om je bloedsuikerspiegel. Denk om je koffie, cola of energiedrank, thee mag natuurlijk ook.

Mocht het echt zo'n dag zijn dat alles tegen zit: zet dan een wekker! Zet na elke pauze een wekker voor de volgende pauze en hou je daar aan!


\section{Op de hoogte houden}
Hou je gebruiker/klant op de hoogte. Afhankelijk van de urgentie van het probleem moet dat meerdere keren per dag of minder. Zorg dat de gebruiker weet waar hij aan toe is. Heb je bijvoorbeeld aan het einde van de dag nog geen oplossing meldt dit dan aan de gebruiker. Vertel hem of haar dat je het overdraagt aan een collega. Geef aan wanneer je er weer verder aan gaat werken. Waar mogelijk betrek de gebruiker in het oplossingsproces, vraag ze eventueel om een mogelijke oplossing te testen.


\chapter{Afronden}
\begin{itemize}
\item Documenteer je oplossing
\item Breng de gebruiker/klant op de hoogte
\item Controleer met de klant/gebruiker of het ook voor de gebruiker/klant opgelost is
\end{itemize}
Als dit allemaal goed gegaan is kan je het ticket in de helpdesk applicatie sluiten.


%%%%%%%%%%%%%%%%%%%%%
%%% Index and End %%%
%%%%%%%%%%%%%%%%%%%%%
\backmatter
\printindex
\end{document}

%%% Last line %%%
